% DISCUSSION - AUTO-UPDATED SECTION
% Last generated: 2026-01-15T07:27:20.548Z

\subsection{Interpretation of Results}

\subsubsection{Quorum Design Implications}

The model is designed to test whether quorum thresholds should be calibrated to expected participation rates rather than set arbitrarily. A hypothesized non-linear relationship between quorum and pass rate implies:

\begin{itemize}
    \item DAOs with low participation should use lower quorums
    \item High-engagement communities can sustain higher thresholds
    \item Dynamic quorum mechanisms may be optimal, adjusting based on recent participation
\end{itemize}

% AUTO-UPDATED: quorum recommendation
Simulation-conditional numeric recommendations are inserted after calibration targets are locked; values here are placeholders until full results are populated.

\subsubsection{Scale Considerations}

A hypothesized power-law decay of participation with scale has profound implications for DAO design:

\begin{enumerate}
    \item Large DAOs must invest in engagement mechanisms
    \item Delegation becomes essential above $\sim$100 members
    \item Sub-DAOs or working groups may preserve participation
\end{enumerate}

This expectation aligns with observations from real DAOs like Uniswap and Compound, though formal calibration is pending.

\subsubsection{Voting Mechanism Trade-offs}

We frame fundamental trade-offs for empirical comparison:

\begin{itemize}
    \item \textbf{Token voting}: Simple, familiar, but plutocratic
    \item \textbf{Quadratic voting}: More egalitarian, but complex and susceptible to sybil attacks
    \item \textbf{Conviction voting}: Favors persistent preferences, but slower decisions
\end{itemize}

No single mechanism dominates; the choice depends on the DAO's values and context.

\subsection{Theoretical Contributions}

This work contributes to theory in several ways:

\begin{enumerate}
    \item \textbf{Quantified trade-offs}: A framework for numerical estimates of previously qualitative relationships
    \item \textbf{Scale effects}: A systematic study design for participation decay in DAO simulations
    \item \textbf{Mechanism comparison}: A controlled comparison setup across voting systems
\end{enumerate}

\subsection{Practical Implications}

For DAO practitioners, this framework suggests:

\begin{enumerate}
    \item \textbf{Start with low quorums}: Begin conservatively and increase if participation supports it
    \item \textbf{Plan for scale}: Design delegation and engagement systems before they're needed
    \item \textbf{Match mechanism to values}: Choose voting systems aligned with community goals
    \item \textbf{Iterate with data}: Use simulation to test changes before on-chain deployment
\end{enumerate}

\noindent
Numeric recommendations should be interpreted as simulation-conditional and comparative rather than universal prescriptions.

\subsection{Connection to Real-World DAOs}

Validation against real DAO pipeline and treasury data is planned as part of the calibration phase. We will report empirical comparisons once relevant on-chain and governance artifacts are integrated into the analysis pipeline.
